\section{Diseño del producto o servicio}
\label{sec:diseno}

El SISTEMA WEB DE INSCRIPCIÓN PARA EL MÓDULO TECNOLÓGICO PRODUCTIVO SAN JULIÁN BTH --- VERSIÓN 2.0 ha sido concebido como una plataforma web integral que abarca la totalidad del ciclo de registro escolar en el MÓDULO TECNOLÓGICO PRODUCTIVO SAN JULIÁN. A diferencia de esquemas parciales, la solución unifica la preinscripción en línea realizada al cierre de la gestión previa con la confirmación de la matrícula formal y la organización académica en la gestión lectiva correspondiente.

\subsection{Arquitectura modular de software}

La solución técnica se estructura bajo una arquitectura cliente-servidor desacoplada en tres capas fundamentales:

\begin{itemize}
    \item \textbf{Capa de Presentación (Frontend):} Desarrollada con estándares abiertos HTML5, hojas de estilo CSS3 con diseño adaptable y scripts de validación reactiva en JavaScript. Se ejecuta directamente en el navegador del usuario (computador o teléfono móvil), permitiendo verificar formatos de carnet y datos obligatorios en tiempo real antes de enviar peticiones a la red.
    \item \textbf{Capa de Lógica del Negocio (Backend):} Implementada en lenguaje PHP 8.x sobre servidor web Apache. Gestiona el procesamiento seguro de solicitudes, la validación de identidad para el restablecimiento de claves, las reglas de asignación documental y la lógica de agrupación en cursos y paralelos.
    \item \textbf{Capa de Persistencia (Base de Datos):} Sustentada en el gestor relacional PostgreSQL 16. Implementa un esquema relacional normalizado con control estricto de transacciones bajo propiedades ACID y bloqueos a nivel de fila (row-level locking), garantizando la consistencia absoluta de las vacantes sin colisiones ni sobrecupos.
\end{itemize}

\begin{figure}[htbp]
    \centering
    \caption{Arquitectura cliente-servidor en tres capas del sistema BTH Versión 2.0.}
    \label{fig:arquitectura_tres_capas}
    \resizebox{0.92\textwidth}{!}{%
        % ============================================================
% Arquitectura Cliente-Servidor en Tres Capas — Cap. 4.1.1
% Proyecto BTH V2.0 — Sistema Web de Inscripción MTP San Julián
% Formato: TikZ vectorial editable (APA 7)
% Uso: % ============================================================
% Arquitectura Cliente-Servidor en Tres Capas — Cap. 4.1.1
% Proyecto BTH V2.0 — Sistema Web de Inscripción MTP San Julián
% Formato: TikZ vectorial editable (APA 7)
% Uso: % ============================================================
% Arquitectura Cliente-Servidor en Tres Capas — Cap. 4.1.1
% Proyecto BTH V2.0 — Sistema Web de Inscripción MTP San Julián
% Formato: TikZ vectorial editable (APA 7)
% Uso: \input{imagenes/arquitectura_tres_capas} dentro de figure
% ============================================================
\begin{tikzpicture}[
    capa/.style={
        rectangle, rounded corners=5pt,
        draw=primary!85, align=center,
        font=\small,
        minimum height=1.7cm,
        inner sep=6pt,
        text width=11.5cm
    },
    flecha/.style={
        -{Stealth[length=3mm,width=2.2mm]},
        thick, draw=primary!85
    },
    flechaInv/.style={
        {Stealth[length=3mm,width=2.2mm]}-,
        thick, draw=primary!60
    },
    etiquetaCapa/.style={
        font=\scriptsize\bfseries, text=primary, align=center
    },
    detalle/.style={
        font=\scriptsize, text=darktext, align=center
    },
    tituloCapa/.style={
        font=\small\bfseries, align=center
    }
]

% --- CAPA 1: PRESENTACIÓN (FRONTEND) ---
\node[capa, fill=lightgray, draw=primary!70] (frontend) at (0, 4.2) {
    \tituloCapa{1. Capa de Presentación \normalfont\textit{(Frontend)}}\\
    \detalle{HTML5 \textbullet{} CSS3 responsivo \textbullet{} JavaScript ES6}\\
    \detalle{\footnotesize Navegador del postulante (PC / móvil) — Validación reactiva en cliente}
};

\node[etiquetaCapa, rotate=90] at (-7.05, 4.2) {CLIENTE};

% Icono usuario sutil (texto)
\node[detalle, font=\tiny\itshape, text=primary!70] at (0, 5.32) {$\sim$400 postulantes/año \textbullet{} 7 especialidades \textbullet{} $<$1.5 MB/trámite};

% --- FLECHA FRONTEND ↔ BACKEND ---
\draw[flecha] (0, 3.30) -- (0, 2.30);
\node[detalle, font=\tiny, fill=white, inner sep=1pt] at (2.15, 2.80) {HTTPS / JSON};
\node[detalle, font=\tiny, fill=white, inner sep=1pt] at (-2.10, 2.80) {Validación};
\draw[flechaInv, dashed, draw=primary!45] (0, 2.30) -- (0, 3.30);

% --- CAPA 2: LÓGICA DE NEGOCIO (BACKEND) ---
\node[capa, fill=secondary!14, draw=primary!70] (backend) at (0, 1.2) {
    \tituloCapa{2. Capa de Lógica del Negocio \normalfont\textit{(Backend)}}\\
    \detalle{PHP 8.2 sobre Apache 2.4}\\
    \detalle{\footnotesize Sesiones, validación C.I.+fecha, reglas matrícula, agrupación paralelos}
};

\node[etiquetaCapa, rotate=90] at (-7.05, 1.2) {SERVIDOR};

% --- FLECHA BACKEND ↔ BD ---
\draw[flecha] (0, 0.30) -- (0, -0.70);
\node[detalle, font=\tiny, fill=white, inner sep=1pt] at (2.35, -0.20) {SQL transaccional};
\node[detalle, font=\tiny, fill=white, inner sep=1pt] at (-2.25, -0.20) {ACID};
\draw[flechaInv, dashed, draw=primary!45] (0, -0.70) -- (0, 0.30);

% --- CAPA 3: PERSISTENCIA (BASE DE DATOS) ---
\node[capa, fill=primary, text=white, draw=primary] (bd) at (0, -1.85) {
    {\small\bfseries 3. Capa de Persistencia \normalfont\textit{(Base de Datos)}}\\
    {\scriptsize PostgreSQL 16.4 — Modelo relacional normalizado}\\
    {\scriptsize\itshape Transacciones ACID + bloqueo pesimista \texttt{FOR UPDATE} — 0 sobrecupos}
};

\node[etiquetaCapa, rotate=90, text=white!90] at (-7.05, -1.85) {DATOS};

% --- Línea lateral vertical decorativa (opcional, sutil) ---
\draw[primary!25, line width=2.5pt, rounded corners=3pt] (-6.55, 5.55) -- (-6.55, -2.75);

% --- Nota de rendimiento ---
\node[detalle, font=\tiny\itshape, text=darktext!65, align=center] at (0, -2.95) {Validación OE4: $>$100 peticiones/min \textbullet{} latencia $<$1.2 s \textbullet{} usabilidad $\geq$90\% \textbullet{} errores $<$5\%};

% --- Borde sutil alrededor (desactivado para APA limpio) ---
% \begin{scope}[on background layer]
%   \node[fit=(frontend)(bd), draw=primary!12, rounded corners=7pt, inner sep=7pt] {};
% \end{scope}

\end{tikzpicture}
 dentro de figure
% ============================================================
\begin{tikzpicture}[
    capa/.style={
        rectangle, rounded corners=5pt,
        draw=primary!85, align=center,
        font=\small,
        minimum height=1.7cm,
        inner sep=6pt,
        text width=11.5cm
    },
    flecha/.style={
        -{Stealth[length=3mm,width=2.2mm]},
        thick, draw=primary!85
    },
    flechaInv/.style={
        {Stealth[length=3mm,width=2.2mm]}-,
        thick, draw=primary!60
    },
    etiquetaCapa/.style={
        font=\scriptsize\bfseries, text=primary, align=center
    },
    detalle/.style={
        font=\scriptsize, text=darktext, align=center
    },
    tituloCapa/.style={
        font=\small\bfseries, align=center
    }
]

% --- CAPA 1: PRESENTACIÓN (FRONTEND) ---
\node[capa, fill=lightgray, draw=primary!70] (frontend) at (0, 4.2) {
    \tituloCapa{1. Capa de Presentación \normalfont\textit{(Frontend)}}\\
    \detalle{HTML5 \textbullet{} CSS3 responsivo \textbullet{} JavaScript ES6}\\
    \detalle{\footnotesize Navegador del postulante (PC / móvil) — Validación reactiva en cliente}
};

\node[etiquetaCapa, rotate=90] at (-7.05, 4.2) {CLIENTE};

% Icono usuario sutil (texto)
\node[detalle, font=\tiny\itshape, text=primary!70] at (0, 5.32) {$\sim$400 postulantes/año \textbullet{} 7 especialidades \textbullet{} $<$1.5 MB/trámite};

% --- FLECHA FRONTEND ↔ BACKEND ---
\draw[flecha] (0, 3.30) -- (0, 2.30);
\node[detalle, font=\tiny, fill=white, inner sep=1pt] at (2.15, 2.80) {HTTPS / JSON};
\node[detalle, font=\tiny, fill=white, inner sep=1pt] at (-2.10, 2.80) {Validación};
\draw[flechaInv, dashed, draw=primary!45] (0, 2.30) -- (0, 3.30);

% --- CAPA 2: LÓGICA DE NEGOCIO (BACKEND) ---
\node[capa, fill=secondary!14, draw=primary!70] (backend) at (0, 1.2) {
    \tituloCapa{2. Capa de Lógica del Negocio \normalfont\textit{(Backend)}}\\
    \detalle{PHP 8.2 sobre Apache 2.4}\\
    \detalle{\footnotesize Sesiones, validación C.I.+fecha, reglas matrícula, agrupación paralelos}
};

\node[etiquetaCapa, rotate=90] at (-7.05, 1.2) {SERVIDOR};

% --- FLECHA BACKEND ↔ BD ---
\draw[flecha] (0, 0.30) -- (0, -0.70);
\node[detalle, font=\tiny, fill=white, inner sep=1pt] at (2.35, -0.20) {SQL transaccional};
\node[detalle, font=\tiny, fill=white, inner sep=1pt] at (-2.25, -0.20) {ACID};
\draw[flechaInv, dashed, draw=primary!45] (0, -0.70) -- (0, 0.30);

% --- CAPA 3: PERSISTENCIA (BASE DE DATOS) ---
\node[capa, fill=primary, text=white, draw=primary] (bd) at (0, -1.85) {
    {\small\bfseries 3. Capa de Persistencia \normalfont\textit{(Base de Datos)}}\\
    {\scriptsize PostgreSQL 16.4 — Modelo relacional normalizado}\\
    {\scriptsize\itshape Transacciones ACID + bloqueo pesimista \texttt{FOR UPDATE} — 0 sobrecupos}
};

\node[etiquetaCapa, rotate=90, text=white!90] at (-7.05, -1.85) {DATOS};

% --- Línea lateral vertical decorativa (opcional, sutil) ---
\draw[primary!25, line width=2.5pt, rounded corners=3pt] (-6.55, 5.55) -- (-6.55, -2.75);

% --- Nota de rendimiento ---
\node[detalle, font=\tiny\itshape, text=darktext!65, align=center] at (0, -2.95) {Validación OE4: $>$100 peticiones/min \textbullet{} latencia $<$1.2 s \textbullet{} usabilidad $\geq$90\% \textbullet{} errores $<$5\%};

% --- Borde sutil alrededor (desactivado para APA limpio) ---
% \begin{scope}[on background layer]
%   \node[fit=(frontend)(bd), draw=primary!12, rounded corners=7pt, inner sep=7pt] {};
% \end{scope}

\end{tikzpicture}
 dentro de figure
% ============================================================
\begin{tikzpicture}[
    capa/.style={
        rectangle, rounded corners=5pt,
        draw=primary!85, align=center,
        font=\small,
        minimum height=1.7cm,
        inner sep=6pt,
        text width=11.5cm
    },
    flecha/.style={
        -{Stealth[length=3mm,width=2.2mm]},
        thick, draw=primary!85
    },
    flechaInv/.style={
        {Stealth[length=3mm,width=2.2mm]}-,
        thick, draw=primary!60
    },
    etiquetaCapa/.style={
        font=\scriptsize\bfseries, text=primary, align=center
    },
    detalle/.style={
        font=\scriptsize, text=darktext, align=center
    },
    tituloCapa/.style={
        font=\small\bfseries, align=center
    }
]

% --- CAPA 1: PRESENTACIÓN (FRONTEND) ---
\node[capa, fill=lightgray, draw=primary!70] (frontend) at (0, 4.2) {
    \tituloCapa{1. Capa de Presentación \normalfont\textit{(Frontend)}}\\
    \detalle{HTML5 \textbullet{} CSS3 responsivo \textbullet{} JavaScript ES6}\\
    \detalle{\footnotesize Navegador del postulante (PC / móvil) — Validación reactiva en cliente}
};

\node[etiquetaCapa, rotate=90] at (-7.05, 4.2) {CLIENTE};

% Icono usuario sutil (texto)
\node[detalle, font=\tiny\itshape, text=primary!70] at (0, 5.32) {$\sim$400 postulantes/año \textbullet{} 7 especialidades \textbullet{} $<$1.5 MB/trámite};

% --- FLECHA FRONTEND ↔ BACKEND ---
\draw[flecha] (0, 3.30) -- (0, 2.30);
\node[detalle, font=\tiny, fill=white, inner sep=1pt] at (2.15, 2.80) {HTTPS / JSON};
\node[detalle, font=\tiny, fill=white, inner sep=1pt] at (-2.10, 2.80) {Validación};
\draw[flechaInv, dashed, draw=primary!45] (0, 2.30) -- (0, 3.30);

% --- CAPA 2: LÓGICA DE NEGOCIO (BACKEND) ---
\node[capa, fill=secondary!14, draw=primary!70] (backend) at (0, 1.2) {
    \tituloCapa{2. Capa de Lógica del Negocio \normalfont\textit{(Backend)}}\\
    \detalle{PHP 8.2 sobre Apache 2.4}\\
    \detalle{\footnotesize Sesiones, validación C.I.+fecha, reglas matrícula, agrupación paralelos}
};

\node[etiquetaCapa, rotate=90] at (-7.05, 1.2) {SERVIDOR};

% --- FLECHA BACKEND ↔ BD ---
\draw[flecha] (0, 0.30) -- (0, -0.70);
\node[detalle, font=\tiny, fill=white, inner sep=1pt] at (2.35, -0.20) {SQL transaccional};
\node[detalle, font=\tiny, fill=white, inner sep=1pt] at (-2.25, -0.20) {ACID};
\draw[flechaInv, dashed, draw=primary!45] (0, -0.70) -- (0, 0.30);

% --- CAPA 3: PERSISTENCIA (BASE DE DATOS) ---
\node[capa, fill=primary, text=white, draw=primary] (bd) at (0, -1.85) {
    {\small\bfseries 3. Capa de Persistencia \normalfont\textit{(Base de Datos)}}\\
    {\scriptsize PostgreSQL 16.4 — Modelo relacional normalizado}\\
    {\scriptsize\itshape Transacciones ACID + bloqueo pesimista \texttt{FOR UPDATE} — 0 sobrecupos}
};

\node[etiquetaCapa, rotate=90, text=white!90] at (-7.05, -1.85) {DATOS};

% --- Línea lateral vertical decorativa (opcional, sutil) ---
\draw[primary!25, line width=2.5pt, rounded corners=3pt] (-6.55, 5.55) -- (-6.55, -2.75);

% --- Nota de rendimiento ---
\node[detalle, font=\tiny\itshape, text=darktext!65, align=center] at (0, -2.95) {Validación OE4: $>$100 peticiones/min \textbullet{} latencia $<$1.2 s \textbullet{} usabilidad $\geq$90\% \textbullet{} errores $<$5\%};

% --- Borde sutil alrededor (desactivado para APA limpio) ---
% \begin{scope}[on background layer]
%   \node[fit=(frontend)(bd), draw=primary!12, rounded corners=7pt, inner sep=7pt] {};
% \end{scope}

\end{tikzpicture}

    }
    \notafigura{Elaboración propia con base en \textcite{tanenbaum2021} y Tabla~\ref{tab:especificaciones_tecnicas}. La arquitectura desacopla presentación (HTML5/CSS3/JS), lógica (PHP 8.2/Apache 2.4) y persistencia (PostgreSQL 16.4 ACID \texttt{FOR UPDATE}).}
\end{figure}

\subsection{Módulos funcionales del sistema}

La plataforma integra los siguientes módulos especializados diseñados para dar respuesta directa a las problemáticas diagnosticadas:

\begin{itemize}
    \item \textbf{Módulo de Formulario Dinámico y Preinscripción:} Interfaz intuitiva y liviana que permite al postulante seleccionar su carrera de interés entre las siete especialidades técnicas habilitadas (Construcción Civil, Contaduría Pública, Electrónica, Gastronomía, Mecánica Automotriz, Sistemas Informáticos y Textiles y Confecciones) e ingresar sus datos personales y de contacto sin recargas innecesarias de página.
    \item \textbf{Módulo de Control Transaccional de Vacantes:} Mecanismo automatizado que consulta y reserva cupos en tiempo real en la base de datos PostgreSQL, impidiendo la duplicidad de registros y asegurando que ninguna plaza sea asignada por encima del límite físico de cada taller.
    \item \textbf{Módulo Autónomo de Recuperación de Credenciales:} Servicio de autoservicio (Self-Service Password Reset) que permite al estudiante recuperar su usuario o restablecer su contraseña de manera autónoma e inmediata mediante la validación de datos personales inequívocos (número de C.I. y fecha de nacimiento), erradicando las filas presenciales y evitando la pérdida del cupo vocacional.
    \item \textbf{Módulo Administrativo de Matriculación Documental:} Herramienta administrativa para la gestión lectiva siguiente que permite a la secretaría del BTH registrar, cotejar y auditar en el sistema la presentación física de los requisitos obligatorios: libretas de calificaciones de 1ro a 3ro de secundaria, cédula de identidad de los tutores, fotocopia de carnet de identidad y certificado de nacimiento, completando el expediente digital formal.
    \item \textbf{Módulo de Asignación a Cursos y Paralelos por Unidad Educativa:} Funcionalidad que distribuye y matricula automáticamente a los postulantes aceptados en cursos y paralelos específicos (tales como 4to A o 4to B) en función de la unidad educativa de procedencia a la que pertenecen, garantizando la articulación horaria y el equilibrio numérico en cada aula y taller.
    \item \textbf{Generador de Comprobantes Digitales Oficiales:} Emisión instantánea de fichas de preinscripción y matrícula en formato PDF descargable e imprimible, provistas de un código alfanumérico único de verificación institucional.
    \item \textbf{Panel de Control y Reportes para Directivos:} Interfaz administrativa que ofrece métricas en tiempo real sobre el estado de vacantes por carrera, porcentajes de matriculación documental completada y exportación de listas consolidadas por paralelo.
\end{itemize}

\begin{figure}[htbp]
    \centering
    \caption{Diagrama de paquetes y dependencias de los siete módulos funcionales del sistema.}
    \label{fig:modulos_paquetes}
    \resizebox{0.98\textwidth}{!}{%
        % ============================================================
% Módulos Funcionales del Sistema — Cap. 4.1.2
% 7 Módulos lógicos — Paquetes UML simplificado (solo lógica)
% Proyecto BTH V2.0 — Sistema Web de Inscripción MTP San Julián
% Formato: TikZ vectorial editable (APA 7)
% Uso: % ============================================================
% Módulos Funcionales del Sistema — Cap. 4.1.2
% 7 Módulos lógicos — Paquetes UML simplificado (solo lógica)
% Proyecto BTH V2.0 — Sistema Web de Inscripción MTP San Julián
% Formato: TikZ vectorial editable (APA 7)
% Uso: % ============================================================
% Módulos Funcionales del Sistema — Cap. 4.1.2
% 7 Módulos lógicos — Paquetes UML simplificado (solo lógica)
% Proyecto BTH V2.0 — Sistema Web de Inscripción MTP San Julián
% Formato: TikZ vectorial editable (APA 7)
% Uso: \input{imagenes/modulos_paquetes} dentro de figure
% ============================================================
\begin{tikzpicture}[
    modulo/.style={
        rectangle, rounded corners=3.5pt,
        draw=primary!75, fill=white,
        text=darktext, align=center,
        font=\scriptsize,
        text width=3.42cm, minimum height=1.35cm,
        inner sep=4pt
    },
    moduloDestacado/.style={
        rectangle, rounded corners=3.5pt,
        draw=primary, fill=primary,
        text=white, align=center,
        font=\scriptsize\bfseries,
        text width=3.42cm, minimum height=1.35cm,
        inner sep=4pt
    },
    moduloLight/.style={
        rectangle, rounded corners=3.5pt,
        draw=primary!55, fill=secondary!11,
        text=darktext, align=center,
        font=\scriptsize,
        text width=3.42cm, minimum height=1.35cm,
        inner sep=4pt
    },
    flecha/.style={
        -{Stealth[length=2.6mm,width=2mm]},
        thick, draw=primary!80
    },
    flechaDashed/.style={
        -{Stealth[length=2.6mm,width=2mm]},
        thick, draw=primary!55, dashed
    },
    etiqueta/.style={
        font=\scriptsize\bfseries, text=primary, align=center
    },
    sistema/.style={
        rectangle, rounded corners=7pt,
        draw=primary!35, thick,
        inner sep=6pt
    }
]

% --- Contenedor Sistema ---
\node[sistema, minimum width=13.2cm, minimum height=6.35cm] (sistema) at (0, 0.15) {};

% Título sistema
\node[etiqueta, fill=white, inner sep=2pt] at (0, 3.55) {SISTEMA WEB BTH V2.0 — MÓDULOS FUNCIONALES};

% --- NIVEL SUPERIOR: Preinscripción (entrada) ---
\node[moduloDestacado] (M1) at (0, 2.45) {\textbf{M1. Formulario}\\Dinámico y\\Preinscripción\\\footnotesize\normalfont\textit{7 especialidades}};

% --- NIVEL MEDIO: Núcleo transaccional + Recovery ---
\node[modulo] (M2) at (-4.15, 0.70) {\textbf{M2. Control}\\Transaccional\\de Vacantes\\\footnotesize\normalfont\textit{ACID row-level}};
\node[moduloLight] (M3) at (0, 0.70) {\textbf{M3. Recuperación}\\Autónoma de\\Credenciales\\\footnotesize\normalfont\textit{C.I.+fecha <$2$ min}};
\node[moduloLight] (M6) at (4.15, 0.70) {\textbf{M6. Generador}\\Comprobantes\\PDF\\\footnotesize\normalfont\textit{código verificación}};

% --- NIVEL INFERIOR: Gestión académica ---
\node[modulo] (M4) at (-4.15, -1.25) {\textbf{M4. Matriculación}\\Documental\\\footnotesize\normalfont\textit{libretas 1ro-3ro,}\\[-2pt]\footnotesize\normalfont\textit{C.I. tutor, carnet, cert.}};
\node[modulo] (M5) at (0, -1.25) {\textbf{M5. Asignación}\\Cursos y Paralelos\\\footnotesize\normalfont\textit{por U.E. procedencia}\\[-2pt]\footnotesize\normalfont\textit{4to A / 4to B}};
\node[modulo] (M7) at (4.15, -1.25) {\textbf{M7. Panel Control}\\y Reportes\\\footnotesize\normalfont\textit{vacantes, \% matrícula}};

% --- Conexiones principales (flujo) ---
% M1 -> M2 (reserva)
\draw[flecha] (M1.south) -- (-2.05, 1.42) -- (M2.north);
\node[font=\tiny, text=primary!80, fill=white, inner sep=1pt] at (-2.35, 1.65) {reserva};

% M1 -> M6 (comprobante)
\draw[flechaDashed] (M1.east) -- (4.85, 1.42) -- (M6.north);
\node[font=\tiny, text=primary!65, fill=white, inner sep=1pt] at (3.15, 1.90) {PDF};

% M1 -> M3 (soporte, recuperación)
\draw[flechaDashed] (M1.south) -- (0, 1.42) -- (M3.north);

% M2 -> M4 (matriculación tras reserva)
\draw[flecha] (M2.south) -- (M4.north);

% M2 -> M5 (asignación tras matrícula)
\draw[flecha] (M2.south) -- (-1.55, -0.20) -- (M5.north);
% M4 -> M5
\draw[flecha] (M4.east) -- (M5.west);
\node[font=\tiny, text=primary!80, fill=white, inner sep=1pt] at (-2.05, -0.85) {exp. ok};

% M5 -> M7 y M6 -> M7 (reportes)
\draw[flechaDashed] (M5.east) -- (M7.west);
\draw[flechaDashed] (M6.south) -- (4.15, -0.05) -- (M7.north);

% M2 -> M7 (vacantes tiempo real)
\draw[flechaDashed] (M2.east) -- (2.10, -0.20) -- (4.15, -0.20) -- (M7.north west);

% --- Leyenda inferior ---
\node[font=\tiny\itshape, text=darktext!65, align=center] at (0, -2.85) {Flujo 5 fases Cap.~4.1.3: (1) Preinscripción → (2) Recovery → (3) Comprobante → (4) Matriculación → (5) Asignación \textbullet{} M7 transversal en tiempo real};
\draw[primary!80, thick] ( -6.0, -2.55) -- (6.0, -2.55);

% --- Indicadores pequeños ---
\node[font=\tiny, text=primary!50, align=left] at (-6.15, 3.15) {\textbf{—} flujo};
\node[font=\tiny, text=primary!50, align=left] at (-6.15, 2.85) {\textbf{- -} consulta};

\end{tikzpicture}
 dentro de figure
% ============================================================
\begin{tikzpicture}[
    modulo/.style={
        rectangle, rounded corners=3.5pt,
        draw=primary!75, fill=white,
        text=darktext, align=center,
        font=\scriptsize,
        text width=3.42cm, minimum height=1.35cm,
        inner sep=4pt
    },
    moduloDestacado/.style={
        rectangle, rounded corners=3.5pt,
        draw=primary, fill=primary,
        text=white, align=center,
        font=\scriptsize\bfseries,
        text width=3.42cm, minimum height=1.35cm,
        inner sep=4pt
    },
    moduloLight/.style={
        rectangle, rounded corners=3.5pt,
        draw=primary!55, fill=secondary!11,
        text=darktext, align=center,
        font=\scriptsize,
        text width=3.42cm, minimum height=1.35cm,
        inner sep=4pt
    },
    flecha/.style={
        -{Stealth[length=2.6mm,width=2mm]},
        thick, draw=primary!80
    },
    flechaDashed/.style={
        -{Stealth[length=2.6mm,width=2mm]},
        thick, draw=primary!55, dashed
    },
    etiqueta/.style={
        font=\scriptsize\bfseries, text=primary, align=center
    },
    sistema/.style={
        rectangle, rounded corners=7pt,
        draw=primary!35, thick,
        inner sep=6pt
    }
]

% --- Contenedor Sistema ---
\node[sistema, minimum width=13.2cm, minimum height=6.35cm] (sistema) at (0, 0.15) {};

% Título sistema
\node[etiqueta, fill=white, inner sep=2pt] at (0, 3.55) {SISTEMA WEB BTH V2.0 — MÓDULOS FUNCIONALES};

% --- NIVEL SUPERIOR: Preinscripción (entrada) ---
\node[moduloDestacado] (M1) at (0, 2.45) {\textbf{M1. Formulario}\\Dinámico y\\Preinscripción\\\footnotesize\normalfont\textit{7 especialidades}};

% --- NIVEL MEDIO: Núcleo transaccional + Recovery ---
\node[modulo] (M2) at (-4.15, 0.70) {\textbf{M2. Control}\\Transaccional\\de Vacantes\\\footnotesize\normalfont\textit{ACID row-level}};
\node[moduloLight] (M3) at (0, 0.70) {\textbf{M3. Recuperación}\\Autónoma de\\Credenciales\\\footnotesize\normalfont\textit{C.I.+fecha <$2$ min}};
\node[moduloLight] (M6) at (4.15, 0.70) {\textbf{M6. Generador}\\Comprobantes\\PDF\\\footnotesize\normalfont\textit{código verificación}};

% --- NIVEL INFERIOR: Gestión académica ---
\node[modulo] (M4) at (-4.15, -1.25) {\textbf{M4. Matriculación}\\Documental\\\footnotesize\normalfont\textit{libretas 1ro-3ro,}\\[-2pt]\footnotesize\normalfont\textit{C.I. tutor, carnet, cert.}};
\node[modulo] (M5) at (0, -1.25) {\textbf{M5. Asignación}\\Cursos y Paralelos\\\footnotesize\normalfont\textit{por U.E. procedencia}\\[-2pt]\footnotesize\normalfont\textit{4to A / 4to B}};
\node[modulo] (M7) at (4.15, -1.25) {\textbf{M7. Panel Control}\\y Reportes\\\footnotesize\normalfont\textit{vacantes, \% matrícula}};

% --- Conexiones principales (flujo) ---
% M1 -> M2 (reserva)
\draw[flecha] (M1.south) -- (-2.05, 1.42) -- (M2.north);
\node[font=\tiny, text=primary!80, fill=white, inner sep=1pt] at (-2.35, 1.65) {reserva};

% M1 -> M6 (comprobante)
\draw[flechaDashed] (M1.east) -- (4.85, 1.42) -- (M6.north);
\node[font=\tiny, text=primary!65, fill=white, inner sep=1pt] at (3.15, 1.90) {PDF};

% M1 -> M3 (soporte, recuperación)
\draw[flechaDashed] (M1.south) -- (0, 1.42) -- (M3.north);

% M2 -> M4 (matriculación tras reserva)
\draw[flecha] (M2.south) -- (M4.north);

% M2 -> M5 (asignación tras matrícula)
\draw[flecha] (M2.south) -- (-1.55, -0.20) -- (M5.north);
% M4 -> M5
\draw[flecha] (M4.east) -- (M5.west);
\node[font=\tiny, text=primary!80, fill=white, inner sep=1pt] at (-2.05, -0.85) {exp. ok};

% M5 -> M7 y M6 -> M7 (reportes)
\draw[flechaDashed] (M5.east) -- (M7.west);
\draw[flechaDashed] (M6.south) -- (4.15, -0.05) -- (M7.north);

% M2 -> M7 (vacantes tiempo real)
\draw[flechaDashed] (M2.east) -- (2.10, -0.20) -- (4.15, -0.20) -- (M7.north west);

% --- Leyenda inferior ---
\node[font=\tiny\itshape, text=darktext!65, align=center] at (0, -2.85) {Flujo 5 fases Cap.~4.1.3: (1) Preinscripción → (2) Recovery → (3) Comprobante → (4) Matriculación → (5) Asignación \textbullet{} M7 transversal en tiempo real};
\draw[primary!80, thick] ( -6.0, -2.55) -- (6.0, -2.55);

% --- Indicadores pequeños ---
\node[font=\tiny, text=primary!50, align=left] at (-6.15, 3.15) {\textbf{—} flujo};
\node[font=\tiny, text=primary!50, align=left] at (-6.15, 2.85) {\textbf{- -} consulta};

\end{tikzpicture}
 dentro de figure
% ============================================================
\begin{tikzpicture}[
    modulo/.style={
        rectangle, rounded corners=3.5pt,
        draw=primary!75, fill=white,
        text=darktext, align=center,
        font=\scriptsize,
        text width=3.42cm, minimum height=1.35cm,
        inner sep=4pt
    },
    moduloDestacado/.style={
        rectangle, rounded corners=3.5pt,
        draw=primary, fill=primary,
        text=white, align=center,
        font=\scriptsize\bfseries,
        text width=3.42cm, minimum height=1.35cm,
        inner sep=4pt
    },
    moduloLight/.style={
        rectangle, rounded corners=3.5pt,
        draw=primary!55, fill=secondary!11,
        text=darktext, align=center,
        font=\scriptsize,
        text width=3.42cm, minimum height=1.35cm,
        inner sep=4pt
    },
    flecha/.style={
        -{Stealth[length=2.6mm,width=2mm]},
        thick, draw=primary!80
    },
    flechaDashed/.style={
        -{Stealth[length=2.6mm,width=2mm]},
        thick, draw=primary!55, dashed
    },
    etiqueta/.style={
        font=\scriptsize\bfseries, text=primary, align=center
    },
    sistema/.style={
        rectangle, rounded corners=7pt,
        draw=primary!35, thick,
        inner sep=6pt
    }
]

% --- Contenedor Sistema ---
\node[sistema, minimum width=13.2cm, minimum height=6.35cm] (sistema) at (0, 0.15) {};

% Título sistema
\node[etiqueta, fill=white, inner sep=2pt] at (0, 3.55) {SISTEMA WEB BTH V2.0 — MÓDULOS FUNCIONALES};

% --- NIVEL SUPERIOR: Preinscripción (entrada) ---
\node[moduloDestacado] (M1) at (0, 2.45) {\textbf{M1. Formulario}\\Dinámico y\\Preinscripción\\\footnotesize\normalfont\textit{7 especialidades}};

% --- NIVEL MEDIO: Núcleo transaccional + Recovery ---
\node[modulo] (M2) at (-4.15, 0.70) {\textbf{M2. Control}\\Transaccional\\de Vacantes\\\footnotesize\normalfont\textit{ACID row-level}};
\node[moduloLight] (M3) at (0, 0.70) {\textbf{M3. Recuperación}\\Autónoma de\\Credenciales\\\footnotesize\normalfont\textit{C.I.+fecha <$2$ min}};
\node[moduloLight] (M6) at (4.15, 0.70) {\textbf{M6. Generador}\\Comprobantes\\PDF\\\footnotesize\normalfont\textit{código verificación}};

% --- NIVEL INFERIOR: Gestión académica ---
\node[modulo] (M4) at (-4.15, -1.25) {\textbf{M4. Matriculación}\\Documental\\\footnotesize\normalfont\textit{libretas 1ro-3ro,}\\[-2pt]\footnotesize\normalfont\textit{C.I. tutor, carnet, cert.}};
\node[modulo] (M5) at (0, -1.25) {\textbf{M5. Asignación}\\Cursos y Paralelos\\\footnotesize\normalfont\textit{por U.E. procedencia}\\[-2pt]\footnotesize\normalfont\textit{4to A / 4to B}};
\node[modulo] (M7) at (4.15, -1.25) {\textbf{M7. Panel Control}\\y Reportes\\\footnotesize\normalfont\textit{vacantes, \% matrícula}};

% --- Conexiones principales (flujo) ---
% M1 -> M2 (reserva)
\draw[flecha] (M1.south) -- (-2.05, 1.42) -- (M2.north);
\node[font=\tiny, text=primary!80, fill=white, inner sep=1pt] at (-2.35, 1.65) {reserva};

% M1 -> M6 (comprobante)
\draw[flechaDashed] (M1.east) -- (4.85, 1.42) -- (M6.north);
\node[font=\tiny, text=primary!65, fill=white, inner sep=1pt] at (3.15, 1.90) {PDF};

% M1 -> M3 (soporte, recuperación)
\draw[flechaDashed] (M1.south) -- (0, 1.42) -- (M3.north);

% M2 -> M4 (matriculación tras reserva)
\draw[flecha] (M2.south) -- (M4.north);

% M2 -> M5 (asignación tras matrícula)
\draw[flecha] (M2.south) -- (-1.55, -0.20) -- (M5.north);
% M4 -> M5
\draw[flecha] (M4.east) -- (M5.west);
\node[font=\tiny, text=primary!80, fill=white, inner sep=1pt] at (-2.05, -0.85) {exp. ok};

% M5 -> M7 y M6 -> M7 (reportes)
\draw[flechaDashed] (M5.east) -- (M7.west);
\draw[flechaDashed] (M6.south) -- (4.15, -0.05) -- (M7.north);

% M2 -> M7 (vacantes tiempo real)
\draw[flechaDashed] (M2.east) -- (2.10, -0.20) -- (4.15, -0.20) -- (M7.north west);

% --- Leyenda inferior ---
\node[font=\tiny\itshape, text=darktext!65, align=center] at (0, -2.85) {Flujo 5 fases Cap.~4.1.3: (1) Preinscripción → (2) Recovery → (3) Comprobante → (4) Matriculación → (5) Asignación \textbullet{} M7 transversal en tiempo real};
\draw[primary!80, thick] ( -6.0, -2.55) -- (6.0, -2.55);

% --- Indicadores pequeños ---
\node[font=\tiny, text=primary!50, align=left] at (-6.15, 3.15) {\textbf{—} flujo};
\node[font=\tiny, text=primary!50, align=left] at (-6.15, 2.85) {\textbf{- -} consulta};

\end{tikzpicture}

    }
    \notafigura{Elaboración propia. M1 preinscripción articula con M2 control transaccional; M3 recuperación autónoma ($<$2 min) y M6 comprobantes PDF operan como soporte; M4 matriculación documental y M5 asignación por U.E. consolidan la gestión académica; M7 panel transversaliza métricas en tiempo real (ver flujo 5 fases Cap.~4.1.3).}
\end{figure}

\subsection{Flujo operativo del ciclo completo de matrícula}

La operación del sistema se articula en cinco fases correlativas:

\begin{enumerate}
    \item \textbf{Fase 1 (Preinscripción y reserva de cupo):} El postulante ingresa desde cualquier dispositivo móvil o PC, completa el formulario validado en cliente y reserva su vacante en la especialidad elegida durante el cierre de la gestión previa.
    \item \textbf{Fase 2 (Soporte y recuperación autónoma):} En caso de olvido de credenciales, el usuario accede al módulo de recuperación, valida su identidad y restablece su acceso al instante sin acudir a ventanillas físicas.
    \item \textbf{Fase 3 (Emisión de ficha preliminar):} El sistema genera y descarga el comprobante digital oficial de preinscripción con el que el estudiante acredita su cupo reservado.
    \item \textbf{Fase 4 (Matriculación documental formal):} Al inicio de la gestión escolar siguiente, el estudiante presenta sus documentos físicos en secretaría, donde el personal verifica y tilda en la plataforma las libretas de 1ro a 3ro de secundaria, cédula de tutores, carnet y certificado de nacimiento.
    \item \textbf{Fase 5 (Asignación académica en cursos y paralelos):} El sistema agrupa a los estudiantes en sus paralelos correspondientes según su colegio de origen, consolidando las listas de clase oficiales del BTH.
\end{enumerate}

\subsection{Especificaciones técnicas del software}

En la Tabla~\ref{tab:especificaciones_tecnicas} se resumen los parámetros técnicos, tecnologías aplicadas y requerimientos de entorno que sustentan la plataforma.

\begin{table}[htbp]
    \centering
    \caption{Especificaciones técnicas de los componentes del sistema web.}
    \label{tab:especificaciones_tecnicas}
    \begin{tabularx}{\textwidth}{L L L}
        \toprule
        \textbf{Componente / Módulo} & \textbf{Tecnología / Parámetro} & \textbf{Especificación / Característica} \\
        \midrule
        Vista (View) & HTML5, CSS3, JavaScript ES6 + plantillas PHP & Interfaz responsiva adaptable a móviles y PC, validación reactiva y renderizado de comprobantes PDF. \\
        Controlador (Controller) & PHP 8.2 sobre Apache 2.4 & Enrutamiento, control de sesiones, validación C.I. + fecha de nacimiento y reglas de matrícula/paralelos. \\
        Modelo (Model) & PHP-PDO + PostgreSQL 16.4 & Clases de dominio y esquema relacional normalizado, transacciones ACID y bloqueo pesimista a nivel de fila (\texttt{FOR UPDATE}). \\
        Autoservicio de Credenciales & Validación de identidad con C.I. y fecha de nacimiento & Restablecimiento autónomo de contraseñas ($<$ 2 min). \\
        Matriculación Documental & Módulo de verificación en base de datos & Registro y auditoría digital de libretas de 1ro a 3ro, cédula de tutores, carnet y certificado. \\
        Asignación Académica & Distribución por cohorte y paralelo (4to A, 4to B) según colegio de procedencia & Agrupación automatizada en secciones. \\
        Comprobantes Digitales & Generador PDF (FPDF/TCPDF) & Emisión instantánea de ficha de inscripción con código alfanumérico único de verificación. \\
        Eficiencia y Conectividad & Optimización y compresión web & Consumo de datos móviles $<$ 1.5 MB por trámite, operable en redes 3G/4G locales. \\
        Validación y Pruebas (OE4) & Caja negra, concurrencia y estrés & $>$100 peticiones/min sin caídas, latencia $<$1.2 s, 0 sobrecupos, usabilidad $\geq$90\% y tasa errores $<$5\%. \\
        \bottomrule
    \end{tabularx}
    \notatabla{Especificaciones técnicas verificadas en el entorno operativo del BTH San Julián.}
\end{table}


\subsection{Utilidad institucional y estándares de calidad}

La utilidad de la plataforma es directa e inmediata: erradica las extenuantes filas nocturnas a la intemperie, descentraliza la atención institucional, democratiza el acceso a las plazas técnicas y otorga total transparencia a la gestión del BTH en San Julián.

Desde la perspectiva de calidad del software (ISO/IEC 25010 y directrices W3C), el sistema destaca por su alta usabilidad, con navegación intuitiva adaptada a usuarios de diversa destreza digital; fiabilidad y seguridad, resguardando la privacidad de los datos personales con contraseñas cifradas y control transaccional ACID; y eficiencia de desempeño, habiendo sido optimizado para consumir un promedio inferior a 1.5 megabytes de datos móviles por trámite, asegurando una operatividad ágil aun bajo coberturas celulares intermitentes en las zonas rurales del municipio.
